\rtmtight
\begin{tblr}{width=\textwidth,colspec={|Q[c,m,2.1cm]|X[l,m]|},hlines,vlines,hspan=minimal,rowsep=1.5pt,colsep=3pt,row{1}={bg=headgray,font=\bfseries}}
\SetCell{c,m}\hfil\logo\hfil & \textbf{POLITEKNIK NEGERI MALANG}\par\textbf{JURUSAN TEKNOLOGI INFORMASI}\par\textbf{PROGRAM STUDI: D4 TEKNIK INFORMATIKA} \\
\SetCell[c=2]{l} \textbf{RENCANA TUGAS MAHASISWA (RTM) DAN RUBRIK PENILAIAN} & \\
Nama Mata Kuliah & Pengembangan Karir \\
Kode Mata Kuliah & RTI256002 \quad \textbf{sks / Jam perminggu:} 2 SKS/4 jam \quad \textbf{Semester:} 6 \\
Dosen Pengampu & \begin{enumerate}\item Agung Nugroho Pramudhita, ST., MT.\item Hendra Pradibta\end{enumerate} \\
\end{tblr}
\assessmentblock{Tugas Analisis dan Refleksi Karir}{Tugas terstruktur -- analisis, refleksi, dan portofolio bertahap}{SCPMK0102-05301, SCPMK0303-05301, SCPMK0303-05302, SCPMK0303-05303}{Mahasiswa menyelesaikan tugas analisis dan refleksi karir secara bertahap setiap minggu, mencakup pemahaman konsep pengembangan karir, identifikasi potensi diri, eksplorasi profesi, strategi personal branding, dan komponen perencanaan karir.}{Membaca dan merangkum materi, merefleksikan konsep terhadap konteks diri sendiri, menjawab pertanyaan analisis, mendokumentasikan evidence belajar, dan mengumpulkan hasil sesuai jadwal.}{Laporan/jurnal refleksi, ringkasan analisis, portofolio digital bertahap, dan/atau dokumen self-assessment.}{Ketepatan dan kedalaman analisis topik karir & 11 & Analisis topik menunjukkan pemahaman mendalam, relevansi dengan konteks profesional informatika, dan dikaitkan dengan potensi serta rencana karir diri sendiri secara spesifik. & Analisis cukup baik dan mencakup topik utama tetapi relevansi dengan konteks diri atau kedalaman argumen masih terbatas. & Analisis tidak menunjukkan pemahaman topik karir yang memadai atau tidak relevan dengan konteks profesional. \\ \\
Refleksi diri dan kesesuaian dengan potensi karir & 10 & Refleksi diri menunjukkan kesadaran mendalam tentang minat, bakat, dan potensi yang dikaitkan secara sistematis dengan peluang karir yang realistis dan terukur. & Refleksi ada dan mencakup aspek utama tetapi belum dikaitkan secara sistematis dengan peluang karir atau masih terlalu umum. & Refleksi tidak bermakna, bersifat deskriptif semata, atau tidak menunjukkan kesadaran diri yang memadai sebagai dasar perencanaan karir. \\ \\
Kebaruan ide dan kreativitas dalam pengembangan karir & 7 & Gagasan baru disampaikan dengan argumentasi yang relevan dan kreatif, disertai contoh konkret dan kaitan dengan perkembangan industri informatika. & Ada ide baru tetapi masih umum, kurang argumentasi, atau tidak dikaitkan dengan konteks industri yang relevan. & Tidak ada kebaruan ide atau gagasan tidak relevan dengan pengembangan karir di bidang informatika. \\}

\assessmentblock{UTS Presentasi Portofolio Karir Awal}{Presentasi individual}{SCPMK0102-05301, SCPMK0303-05301, SCPMK0303-05302}{Mahasiswa mempresentasikan portofolio karir awal yang disusun secara sistematis, mencakup profil diri, analisis minat dan potensi, profil karir yang dipilih, strategi personal branding, dan rencana awal pengembangan karir.}{Menyusun portofolio karir awal secara individual, mempresentasikan secara lisan dengan bahan visual yang profesional, dan mempertahankan argumentasi rencana karir berdasarkan analisis diri.}{Bahan presentasi (slide), portofolio karir awal tertulis, dan kemampuan komunikasi lisan saat sesi tanya jawab.}{Kelengkapan dan kualitas portofolio karir awal & 9 & Portofolio mencakup profil diri, analisis minat/potensi, profil karir yang dipilih, dan strategi personal branding yang terhubung secara sistematis dan menunjukkan pemikiran yang matang. & Portofolio mencakup komponen utama tetapi hubungan antar komponen atau kedalaman analisis masih terbatas atau belum sistematis. & Portofolio tidak lengkap, tidak menunjukkan analisis karir yang bermakna, atau tidak dapat digunakan sebagai dasar perencanaan. \\ \\
Kemampuan presentasi dan komunikasi profesional & 8 & Presentasi disampaikan dengan percaya diri, terstruktur, menggunakan bahasa profesional yang tepat, dan bahan visual yang mendukung argumen secara efektif. & Presentasi cukup jelas dan terstruktur tetapi masih ada kelemahan dalam kepercayaan diri, konsistensi bahasa profesional, atau kualitas bahan visual. & Presentasi tidak terstruktur, tidak menggunakan bahasa profesional, atau tidak menunjukkan kemampuan komunikasi yang memadai. \\ \\
Argumentasi rencana karir berdasarkan analisis diri & 8 & Rencana karir diargumentasikan dengan logis berdasarkan analisis kekuatan, kelemahan, peluang, dan risiko yang teridentifikasi secara konkret dan dapat dipertahankan. & Argumentasi ada tetapi koneksi dengan analisis diri masih lemah atau sebagian argumen belum dapat dipertahankan saat tanya jawab. & Tidak ada argumentasi yang bermakna, rencana tidak berdasarkan analisis diri, atau tidak dapat merespons pertanyaan dengan baik. \\}

\assessmentblock{PBL Career Development Plan}{Project Based Learning -- Penyusunan dan Presentasi Proyek}{SCPMK0303-05303, SCPMK0303-05302, SCPMK0102-05301}{Mahasiswa menyusun dan mempresentasikan Career Development Plan (CDP) yang komprehensif sebagai proyek akhir individual, mencakup tujuan karir SMART, analisis gap kompetensi, strategi pengembangan diri, timeline pencapaian, dan rencana evaluasi berkala.}{Menyusun CDP secara sistematis berdasarkan seluruh pembelajaran sepanjang semester, mempresentasikan hasilnya secara profesional di hadapan dosen dan rekan sejawat, serta mempertahankan keputusan dan strategi yang dipilih berdasarkan analisis yang kuat.}{Dokumen CDP tertulis yang lengkap, bahan presentasi profesional, dan kemampuan mempertahankan argumen saat sesi tanya jawab dan peer review.}{Kelengkapan dan sistematika Career Development Plan & 15 & CDP mencakup tujuan karir SMART, analisis gap kompetensi yang mendalam, roadmap pengembangan, timeline yang realistis, strategi pencapaian, dan rencana evaluasi yang tersusun secara sistematis dan koheren. & CDP mencakup komponen utama tetapi timeline, detail strategi, atau analisis gap kompetensi masih kurang lengkap atau belum cukup spesifik. & CDP tidak lengkap, tidak sistematis, atau tidak dapat digunakan sebagai panduan pengembangan karir yang nyata. \\ \\
Kedalaman analisis dan realisme rencana karir & 14 & Rencana karir realistis, terukur, dan didasarkan pada analisis mendalam tentang kompetensi diri, peluang industri informatika, tren karir, dan potensi tantangan yang diantisipasi. & Rencana cukup realistis tetapi analisis kompetensi, peluang industri, atau antisipasi tantangan belum cukup mendalam. & Rencana tidak realistis, tidak terukur, atau tidak berdasarkan analisis yang memadai tentang kompetensi dan peluang karir. \\ \\
Kualitas presentasi dan kemampuan mempertahankan CDP & 11 & CDP dipresentasikan secara profesional dengan bahan visual yang kuat, argumen dipertahankan dengan evidence konkret, dan mampu menjawab pertanyaan dengan baik serta terbuka terhadap masukan. & Presentasi cukup baik dan argumen dapat disampaikan, tetapi kemampuan mempertahankan detail CDP atau merespons pertanyaan yang mendalam masih terbatas. & Tidak dapat mempertahankan CDP secara bermakna, presentasi tidak profesional, atau tidak dapat merespons pertanyaan dengan argumen yang valid. \\}

\assessmentblock{UAS Evaluasi Akhir Pengembangan Karir}{Evaluasi reflektif dan presentasi singkat}{SCPMK0303-05301, SCPMK0303-05302}{Mahasiswa melakukan evaluasi akhir atas proses pengembangan karir selama semester, merefleksikan perkembangan diri, dan menunjukkan konsistensi nilai etika profesional dalam seluruh portofolio dan keputusan karir yang telah dibuat.}{Menyusun refleksi akhir secara tertulis dan/atau lisan yang merangkum proses belajar, perkembangan diri, langkah konkret ke depan, dan demonstrasi nilai etika profesional dalam konteks karir.}{Laporan refleksi akhir, ringkasan evaluasi diri, dan/atau presentasi singkat.}{Kualitas refleksi akhir dan sintesis pembelajaran karir & 4 & Mahasiswa mampu merefleksikan seluruh proses pembelajaran secara mendalam, mengidentifikasi perkembangan diri yang konkret, dan merumuskan langkah selanjutnya yang realistis dan terukur. & Refleksi ada dan mencakup aspek utama pembelajaran tetapi sintesis antar topik atau rencana tindak lanjut masih terbatas dan belum spesifik. & Tidak dapat merumuskan refleksi yang bermakna tentang perkembangan karir atau tidak menunjukkan kemampuan sintesis pembelajaran. \\ \\
Konsistensi nilai etika profesional dalam karir & 3 & Komitmen terhadap etika profesional dan integritas karir ditunjukkan secara konsisten melalui seluruh portofolio, keputusan, dan presentasi yang dibuat sepanjang semester. & Nilai etika disebutkan dan sebagian tercermin dalam karya atau presentasi tetapi belum konsisten atau belum terintegrasi dalam setiap keputusan karir. & Tidak menunjukkan pemahaman nilai etika profesional atau nilai tersebut tidak tercermin dalam portofolio dan keputusan karir. \\}

\begin{tblr}{width=\textwidth,colspec={|Q[l,m,3.2cm]|X[l,m]|},hlines,vlines,hspan=minimal,row{1-3}={bg=headgray},rowsep=1.5pt,colsep=3pt}
Jadwal Pelaksanaan: & Minggu 1--17 sesuai RPS. Setiap evaluasi memiliki RTM dan rubrik multi-kriteria. \\
Lain-Lain Yang Diperlukan: & LMS, platform presentasi daring, alat self-assessment karir, dan perangkat presentasi. \\
Pustaka: & Mengacu pustaka utama dan pendukung pada bagian RPS. \\
\end{tblr}
